\subsection{IMAGE OF A CONNECTED SET UNDER A CONTINUOUS MAPPING}

PROPOSITION 4. Let A be a connected subset of a topological space X, and let f be a continuous mapping of X into a topological space X'. Then \(f(\mathbf{A})\) is connected.

Suppose \(f(\mathbf{A})\) is not connected. Then there exist two sets \(\mathbf{M}'\) , \(\mathbf{N}'\) which are open in \(f(\mathbf{A})\) and which form a partition of \(f(\mathbf{A})\) ; hence \(\mathbf{A} \cap \overline{f}^{1}(\mathbf{M}')\) and \(\mathbf{A} \cap \overline{f}^{1}(\mathbf{N}')\) are open in \(\mathbf{A}\) and form a partition of \(\mathbf{A}\) ; this contradicts the hypothesis that \(\mathbf{A}\) is connected.

The inverse image of a connected set under a continuous mapping need not be connected; consider for example a mapping of a discrete space into a space consisting of one point.

From Proposition 4 we derive another characterization of non-connected spaces:

PROPOSITION 5. For a topological space X to be not connected it is necessary and sufficient that there exists a surjective continuous mapping of X onto a discrete space containing more than one point.

The condition is sufficient by Proposition 4. Conversely, if X is not connected, there exist two non-empty disjoint open subsets A, B whose union is X, and the mapping f of X onto a discrete space of two elements \(\{a, b\}\) , defined by \(f(A) = \{a\}\) and \(f(B) = \{b\}\) , is continuous.

